\begin{conjecture}[Dirichlet 挠点上的 \texorpdfstring{$\mathrm{RH}(u)$}{RH(u)} 失败应在单位圆上稠密]\label{conj:xi-dirichlet-torsion-rhu-failure-dense}
设
\[
\mathcal E_{\mathrm{tor}}
:=
\{u\in \mathrm{Tor}(\TT)\setminus\{1\}:\ \text{对应扭曲上的 RH 型平方根律失败}\}.
\]
则应有如下更强现象：除有限多个例外外，全部非平凡挠点都属于 \(\mathcal E_{\mathrm{tor}}\)。特别地，
\[
\overline{\mathcal E_{\mathrm{tor}}}=\TT.
\]
\end{conjecture}
\begin{remark}
该猜想与本文已有三条证据完全同向：其一，定理 \ref{thm:xi-output-potential-large-odd-primes-bad} 已给出充分大奇素数模数全部为坏模数；其二，推论 \ref{cor:xi-output-potential-bad-moduli-density-one} 已证明坏模数集合在 \(\NN\) 中具有自然密度 \(1\)；其三，结论 \ref{con:xi-time-part9l-dirichlet-twist-divisibility-monotonicity} 表明平方根门槛失效会沿整除链向上传播。因而，现有结果已经把扭曲平方根律的失败推进到一个自然密度为 \(1\) 的模数族上；本猜想只是把这种“模数典型性”进一步提升为“挠点典型性”，并断言其在单位圆上形成稠密失败集。
\end{remark}

